\subsection*{Что измеряет парный расчёт Монте-Карло}

В объёмно-зависимом кластерном разложении C46 на каждой конфигурации
заполнения \(\sigma\) и удельном объёме \(v\) известны суммы по орбитам
\(S_\gamma(\sigma)\) для всей ячейки. Полная эффективная энергия
\(H_m(\sigma,v)=\sum_\gamma S_\gamma J_{\gamma,m}(v,T)\) включает
статическую часть MTP и колебательную свободную энергию для выбранной
статистики \(m\). На девяти узлах по объёму подгоняются все 20
коэффициентов, между узлами используется кубический сплайн.
В Монте-Карло меняются химические типы атомов и общий объём
при неизменной гексагональной форме. Внешнее давление равно нулю.

Чтобы измерить производную свободной энергии по составу \(x=N_{\rm Al}/N\),
в окне \(l\) к \(H_m\) добавляли известное смещение
\(-N\mu_0 x+Nk_{\rm B}T\kappa(x-x_{0,l})^2\).
Для обеих моделей использованы \(\kappa=100\) и
\(\mu_0=3{,}37431251\)~эВ на замену Ti на Al. По средним составам
определяли химическое поле и общую касательную,
\[
 \widehat{\Delta\mu}_{m,l}
 =\mu_0-2k_{\rm B}T\kappa(\langle x\rangle_{m,l}-x_{0,l}),\qquad
 \int_{x_\alpha}^{x_{\alpha_2}}
 [\widehat{\Delta\mu}_{m}(x)-\Delta\mu^*_m]\,dx=0 .
\]
Поля на концах интегрирования равны \(\Delta\mu_m^*\). Изотерму между
окнами интерполировали кусочно-кубической функцией с сохранением формы
(PCHIP). Положение двух контактов зависит от дискретизации изотермы,
поэтому основной контроль её разрешения выполнен удвоением плотности окон.

В 13 общих окнах центры \(x_0\) менялись от 0,13 до 0,25 с шагом
0,01. Сохраняли состав, объём и все суммы по орбитам. Это позволяет
точно вычислить \(\Delta H=H_{\rm cl}-H_{\rm q}\) на каждом кадре,
включая константу и точечную орбиту в текущем объёме.
Для данной ячейки и окна химическое смещение одинаково в обеих моделях.
Отношение весов классической и квантовой моделей пропорционально
\(\exp[-\Delta H/(k_{\rm B}T)]\), где в экспоненте стоит энергия
всей ячейки. Двусторонняя оценка Беннетта объединяет траектории
обеих моделей в одном окне. Благодаря общим кадрам оценка разности
границ сохраняет корреляцию между двумя рассчитанными изотермами.

Для \(N=5120\) в каждом из 13 окон одна серия начиналась в
упорядоченной конфигурации с классическими коэффициентами, другая~---
в случайной конфигурации с квантовыми. На серию приходилось 5000
проходов уравновешивания и 20\,000 производственных проходов, с
записью каждого десятого. В дополнительные 12 окон с центрами от
0,135 до 0,245 с шагом 0,01 запускали классические случайные старты
той же длины. В этих окнах квантовые средние получены обратным
перевзвешиванием с множителем \(\exp[+\Delta H/(k_{\rm B}T)]\).
В каждой бутстрэп-реализации для обеих моделей брали одни и те же
временные блоки из каждой траектории и заново находили обе общие
касательные. Интервал разности получали по распределению парных
разностей составов. Чувствительность к числу узлов, размеру ячейки
и выбору интерполяции оценивается отдельно.

\paragraph{Проверка перекрытия и начальных состояний.}
В 13 исходных окнах \(N=5120\) минимальная доля эффективных кадров
Киша после перевзвешивания составила \(0{,}904\); она измеряет
концентрацию весов до поправки на временную корреляцию. Число
временных блоков менялось от 19 до 200 на отдельную траекторию.
По квантовым случайным и классическим упорядоченным стартам раздельно
сдвиг \(x_\alpha\) равен \(-0{,}039\) и \(-0{,}057\)~п.п.,
а сдвиг \(x_{\alpha_2}\) равен \(+0{,}006\) и \(-0{,}012\)~п.п.
Абсолютный правый контакт классической модели по этим двум сериям
равен \(22{,}020\) и \(21{,}840\) ат.\% Al. Разность \(0{,}180\)~п.п.
имеет условный 95\%-й интервал \([0{,}049;\,0{,}294]\)~п.п. и составляет
2,87 суммарной стандартной ошибки. Для разностей парных поправок
между этими сериями 95\%-е интервалы включают нуль на обеих границах.
Увеличение длины блоков вдвое сохранило выводы по парным сдвигам.
Для 13 узлов обе границы разрешены во всех 768 совместных
бутстрэп-реализациях; при одновременной проверке 13 и 25 окон~---
во всех 512 реализациях. Разность полей для густой сетки составляет
\(-1{,}928\)~мэВ на замену.
